\begin{table}[t]
\caption{Total Effect Without Topic Adjustment — H2 (Probabilistic Time)} 
\fontsize{12.0pt}{14.0pt}\selectfont
\begin{tabular*}{\linewidth}{@{\extracolsep{\fill}}lccc}
\toprule
  & (1) Full Sample† & (2) Total† & (3) Main† \\ 
\midrule\addlinespace[2.5pt]
Invasion & -0.355*** & -0.382*** & -0.347*** \\ 
 & (0.039) & (0.042) & (0.027) \\ 
State-Owned x Treat & 0.17 & 0.174 & 0.222* \\ 
{} & {(0.08)} & {(0.088)} & {(0.069)} \\ 
R2 & 0.039 & 0.020 & 0.024 \\ 
R2 Adj. & 0.039 & 0.020 & 0.023 \\ 
Topic FEs & X & - & X \\ 
Military Included & X & X & - \\ 
Source Z-Score & X & X & X \\ 
N (within bandwidth) & 38,264 & 38,218 & 31,405 \\ 
Clusters & 43 & 43 & 43 \\ 
\bottomrule
\end{tabular*}
\begin{minipage}{\linewidth}
\vspace{.05em}
+ p < 0.1, * p < 0.05, ** p < 0.01, *** p < 0.001\\
The total-effect model removes the topic adjustment. It is a distinct estimand,
not a robustness check: it lets the topic composition of coverage move rather
than holding it fixed. Outcome, source clustering and bandwidth procedure are
otherwise unchanged, with the bandwidth re-selected without topic. Model 1,
the full sample model, is unreported in the main text. It is the hypothesis
2 interaction estimated on the non-NATO sample, but with all topics included.
Model 3 is Table 2 Model 5, the main H2 specification. The main specification
omits military coverage, so the total effect is estimated on the full sample
and Model 1 is its like-for-like comparator. Model 1 against Model 2 isolates
the topic adjustment; Model 2 against Model 3 shows the effect of the sample
restriction. † Standard errors computed via wild cluster restricted (WCR)
bootstrap (Cameron, Gelbach \& Miller 2008) with Webb six-point weights (B=4999,
null imposed), clustered by source domain.\\
\end{minipage}
\end{table}

